\documentclass[11pt, letterpaper]{article}

\usepackage[margin=1in]{geometry}
\usepackage{amsmath}
\usepackage{amssymb}
\usepackage{graphicx}
\usepackage[utf8]{inputenc}
\usepackage{hyperref}
\usepackage{array}

\setlength{\parindent}{0pt}
\setlength{\parskip}{0.8em}
\renewcommand{\labelitemi}{\textbullet}
\renewcommand{\labelitemii}{$\circ$}

\hypersetup{
    colorlinks=true,
    linkcolor=blue,
    filecolor=magenta,
    urlcolor=cyan,
}

\newcommand{\question}[2]{
    \subsection*{Question #1: #2}
}
\newcommand{\points}[1]{
    \textbf{(#1 Points)}
}

\begin{document}

\begin{center}
    {\Huge \textbf{Assignment 1 Grading Rubric}} \\
    \vspace{0.2cm}
    {\large IEDA 1250: Optimization for Decision Making (Summer, 2026)} \\
    \vspace{0.5cm}
    \hrule
\end{center}

\question{1}{Linear Programming -- Graphical Solutions \points{10}}

\begin{itemize}
    \item Correctly analyzes LP (a), including the graphical feasible region or
    equivalent reasoning, and concludes that the LP is unbounded. \points{2.5}

    \item Correctly analyzes LP (b), including the graphical feasible region or
    equivalent reasoning, and concludes that the LP is unbounded. \points{2.5}

    \item Correctly analyzes LP (c), identifies a unique optimal solution
    \[
    x_1^*=1,\qquad x_2^*=0,
    \]
    and reports the corresponding objective value. \points{2.5}

    \item Correctly analyzes LP (d), identifies infinitely many optimal solutions on
    \[
    x_1+3x_2=3
    \]
    between
    \[
    \left(\frac{3}{2},\frac{1}{2}\right)
    \quad \text{and} \quad
    (3,0),
    \]
    and reports the corresponding objective value. \points{2.5}
\end{itemize}


\question{2}{Duality -- Transforming a General LP to Standard Form \points{10}}

\begin{itemize}
    \item Correctly substitutes
    \[
    x_1=-\bar{x}_1,\qquad x_2=\bar{x}_2-\underline{x}_2,
    \]
    with all new variables nonnegative, and rewrites the objective and constraints
    in the new variables. \points{2}

    \item Correctly replaces the equality constraint with one ``$\leq$'' constraint
    and one ``$\geq$'' constraint. \points{2}

    \item Correctly converts all ``$\geq$'' inequalities to ``$\leq$'' inequalities
    so that the primal is in a form with only ``$\leq$'' constraints and
    nonnegative variables. \points{2}

    \item Correctly writes the dual of the transformed LP, including the correct
    objective, constraints, and nonnegativity restrictions on the dual variables.
    \points{2}

    \item Correctly simplifies the dual LP to the equivalent compact form with one
    free dual variable and one nonnegative dual variable. \points{2}
\end{itemize}


\question{3}{Complementary Slackness \points{10}}

\begin{itemize}
    \item Correctly writes or uses the dual LP:
    \[
    \max \quad 4y_1+3y_2
    \]
    with the five dual constraints corresponding to $x_1,\ldots,x_5$ and
    $y_1,y_2\geq 0$. \points{2}

    \item Correctly evaluates the dual constraints at
    \[
    y_1^*=0.8,\qquad y_2^*=0.6,
    \]
    and identifies which dual constraints are tight and which have slack. \points{2}

    \item Correctly applies complementary slackness to conclude
    \[
    x_2^*=x_3^*=x_4^*=0.
    \]
    \points{2}

    \item Correctly uses $y_1^*>0$ and $y_2^*>0$ to make both primal constraints
    tight and solves
    \[
    x_1+3x_5=4,\qquad 2x_1+x_5=3.
    \]
    \points{2}

    \item Correctly reports the optimal primal solution and objective value:
    \[
    x_1^*=1,\quad x_2^*=0,\quad x_3^*=0,\quad x_4^*=0,\quad x_5^*=1,
    \qquad z^*=5.
    \]
    \points{2}
\end{itemize}


\question{4}{Linear Programming Model and Excel Solution -- PaGas Water Operation \points{10}}

\begin{itemize}
    \item Correctly defines and uses the five decision variables
    $X_{AF},X_{BF},X_{AR},X_{BR},X_F$. \points{1}

    \item Correctly formulates the total cost objective:
    \[
    \min \quad 200X_{AF}+100X_{AR}+180X_{BF}+120X_{BR}+70X_F.
    \]
    \points{2}

    \item Correctly includes the water balance constraints at Plant A, Plant B, and
    the fracking site:
    \[
    X_{AF}+X_{AR}=20,\quad
    X_{BF}+X_{BR}=25,\quad
    X_{AF}+X_{BF}+X_F=50.
    \]
    \points{3}

    \item Correctly imposes nonnegativity on all decision variables. \points{1}

    \item Provides a valid Excel Solver setup or screenshot showing the model and
    optimal solution. \points{1}

    \item Correctly reports the optimal solution:
    \[
    X_{AF}^*=0,\quad X_{AR}^*=20,\quad X_{BF}^*=25,\quad
    X_{BR}^*=0,\quad X_F^*=25.
    \]
    \points{2}
\end{itemize}


\question{5}{Binary Integer Programming -- Crew Assignment \points{20}}

\begin{itemize}
    \item Correctly defines twelve binary variables
    \[
    x_i=
    \begin{cases}
    1, & \text{if sequence }i\text{ is selected,}\\
    0, & \text{otherwise,}
    \end{cases}
    \qquad i=1,\ldots,12.
    \]
    \points{2}

    \item Correctly formulates the cost-minimization objective:
    \[
    \min \; 2x_1+3x_2+4x_3+6x_4+7x_5+5x_6
    +7x_7+8x_8+9x_9+9x_{10}+8x_{11}+9x_{12}.
    \]
    \points{3}

    \item Correctly writes coverage constraints requiring each flight to be covered
    by at least one selected sequence. Award partial credit for each correctly
    modeled flight coverage constraint. \points{7}

    \item Correctly imposes the exactly-three-crews constraint:
    \[
    \sum_{i=1}^{12}x_i=3.
    \]
    \points{2}

    \item Correctly states the binary restrictions $x_i\in\{0,1\}$ for all
    $i=1,\ldots,12$. \points{2}

    \item Provides a valid Excel Solver setup or screenshot showing the integer
    programming model and optimal solution. \points{2}

    \item Correctly reports one optimal solution and its total cost. Valid optimal
    solutions include
    \[
    x_3=x_4=x_{11}=1
    \]
    or
    \[
    x_1=x_5=x_{12}=1,
    \]
    with all other variables equal to 0 and total cost $18{,}000$. \points{2}
\end{itemize}


\question{6}{Mixed Integer Reformulation of Logical Constraints \points{10}}

\begin{itemize}
    \item Correctly introduces binary variables to indicate whether each logical
    inequality is enforced. \points{2}

    \item Correctly reformulates part (a), ``at least one of two inequalities
    holds,'' using a valid big-$M$ formulation, for example:
    \[
    3x_1+x_2+x_3+x_4 \leq 12+M(1-y_1),
    \]
    \[
    x_1+x_2+x_3+x_4 \leq 15+M(1-y_2),
    \]
    \[
    y_1+y_2\geq 1,\qquad y_1,y_2\in\{0,1\}.
    \]
    \points{4}

    \item Correctly reformulates part (b), ``at least two of three inequalities
    hold,'' using a valid big-$M$ formulation, for example:
    \[
    2x_1+5x_2+x_3+x_4 \leq 30+M(1-y_1),
    \]
    \[
    x_1+3x_2+5x_3+x_4 \leq 40+M(1-y_2),
    \]
    \[
    3x_1+x_2+3x_3+x_4 \leq 60+M(1-y_3),
    \]
    \[
    y_1+y_2+y_3\geq 2,\qquad y_1,y_2,y_3\in\{0,1\}.
    \]
    \points{4}
\end{itemize}


\question{7}{Mixed Binary Integer Programming -- Product Mix \points{20}}

\begin{itemize}
    \item Correctly defines continuous production variables $x_i\geq 0$ and binary
    setup variables $y_i\in\{0,1\}$ for products $i=1,\ldots,4$. \points{3}

    \item Correctly formulates the profit-maximization objective including marginal
    revenues and start-up costs:
    \[
    \max \; 70x_1+60x_2+90x_3+80x_4
    -50000y_1-40000y_2-70000y_3-60000y_4.
    \]
    \points{3}

    \item Correctly imposes the start-up linking constraints
    \[
    0\leq x_i\leq My_i,\qquad i=1,\ldots,4.
    \]
    \points{3}

    \item Correctly models that no more than two products can be produced:
    \[
    y_1+y_2+y_3+y_4\leq 2.
    \]
    \points{2}

    \item Correctly models that Product 3 or Product 4 can be produced only if
    Product 1 or Product 2 is produced:
    \[
    y_3\leq y_1+y_2,\qquad y_4\leq y_1+y_2.
    \]
    \points{3}

    \item Correctly models that exactly one of the two resource constraints holds
    using a valid big-$M$ formulation and an auxiliary binary variable. \points{3}

    \item Provides a valid Excel Solver setup or screenshot and reports the optimal
    solution
    \[
    x_2^*=2000,\qquad y_2^*=1,
    \]
    with all other production and setup variables equal to 0, and maximum total
    profit $80{,}000$. \points{3}
\end{itemize}


\question{8}{Expected Value of a Discrete Random Variable \points{10}}

\begin{itemize}
    \item Correctly computes
    \[
    E[X]=0(0.2)+1(0.5)+2(0.3)=1.1.
    \]
    \points{4}

    \item Correctly computes
    \[
    E[X^2]=0^2(0.2)+1^2(0.5)+2^2(0.3)=1.7.
    \]
    \points{2}

    \item Correctly uses linearity of expectation:
    \[
    E[2X^2+3]=2E[X^2]+3=6.4.
    \]
    \points{4}
\end{itemize}

\end{document}
